\section{From point probabilities to interval probabilities}

The passage from discrete random variables to continuous random variables is one
of the most important conceptual changes in elementary probability. In the
discrete case, probability can be assigned directly to separate values:
\[
P(X=x).
\]
For example, if \(X\) is the result of one fair die roll, then
\[
P(X=4)=\frac16.
\]
The value \(4\) carries a positive amount of probability mass.

For many real measurements this picture is no longer appropriate. If \(X\) is a
waiting time, a height, a temperature, a measurement error, or a lifetime, then
there are not just a few separated possible values. Between any two possible
values there are many other possible values. The natural questions are then not
usually point questions such as
\[
P(X=2.0000000000),
\]
but interval questions such as
\[
P(1.9\leq X\leq 2.1),
\]
or
\[
P(X>5).
\]

The purpose of this chapter is to build the correct language for such questions.
The main idea is simple to state but easy to misunderstand:

\begin{summarybox}
For a continuous random variable, probability is not concentrated at individual
points. Probabilities of intervals are the basic objects.
\end{summarybox}

\subsection*{A first example: choosing a point from an interval}

Imagine that a point is chosen uniformly from the interval \([0,1]\). Let
\[
X=\text{the chosen point}.
\]
The model says that all subintervals of the same length should have the same
probability. Therefore the interval \([0.2,0.5]\), whose length is \(0.3\), should
have probability \(0.3\):
\[
P(0.2\leq X\leq 0.5)=0.3.
\]
Likewise,
\[
P(0.7\leq X\leq 0.9)=0.2.
\]

Now ask for the probability of one exact point, for instance \(X=0.4\). If that
point had some positive probability, say \(c>0\), then every other point should
also have the same probability by uniformity. But the interval \([0,1]\) contains
infinitely many points. Assigning the same positive mass to every point would
make the total probability exceed \(1\). The only reasonable value is
\[
P(X=0.4)=0.
\]

This does not mean that the value \(0.4\) is impossible in the ordinary verbal
sense. The model allows \(0.4\) as a possible value. It means that a single point
has no length and therefore receives no probability mass in this continuous
model.

\begin{remarkbox}
Probability zero does not always mean logical impossibility. In continuous
models, individual points may have probability zero even though they belong to
the set of possible values.
\end{remarkbox}

\subsection*{Why intervals replace points}

In the uniform model on \([0,1]\), probability is proportional to length:
\[
P(a\leq X\leq b)=b-a,
\]
whenever \(0\leq a\leq b\leq 1\). Thus
\[
P(0.1\leq X\leq 0.4)=0.3,
\]
while
\[
P(0.100\leq X\leq 0.400)=0.300.
\]
The number of decimal places used to write the endpoints is irrelevant. What
matters is the interval itself.

A single point is an interval whose length has collapsed to zero:
\[
\{a\}=[a,a].
\]
Therefore, in this model,
\[
P(X=a)=P(a\leq X\leq a)=a-a=0.
\]
This is the first fundamental difference between the discrete and continuous
cases.

\subsection*{Comparison with a discrete model}

Suppose \(Y\) is the result of one fair die roll. Then the possible values are
\[
1,2,3,4,5,6.
\]
The probability law is determined by the point probabilities
\[
P(Y=k)=\frac16,\qquad k=1,2,3,4,5,6.
\]
An event such as \(Y\in\{2,3,4\}\) is handled by adding point masses:
\[
P(Y\in\{2,3,4\})=P(Y=2)+P(Y=3)+P(Y=4)=\frac36.
\]

For the uniform random point \(X\) on \([0,1]\), the corresponding calculation is
not a sum over points. The event \(0.2\leq X\leq 0.5\) contains infinitely many
points, but each individual point has probability zero. The event has positive
probability because the whole interval has positive length:
\[
P(0.2\leq X\leq 0.5)=0.5-0.2=0.3.
\]

\begin{center}
\begin{tabular}{lll}
\toprule
Situation & Basic object & Probability is computed by \\
\midrule
Discrete model & Points & Adding probability masses \\
Continuous model & Intervals & Measuring accumulated density \\
\bottomrule
\end{tabular}
\end{center}

The phrase ``accumulated density'' will be made precise in the next sections.
For now, the key idea is that a continuous probability law spreads probability
across a continuum of values.

\subsection*{Endpoint conventions}

In discrete probability, the difference between
\[
P(Y<4)
\]
and
\[
P(Y\leq 4)
\]
can be important, because the value \(4\) may have positive probability.

In continuous models with no point masses, endpoints do not change interval
probabilities. For example, if \(X\) is uniformly distributed on \([0,1]\), then
\[
P(0.2<X<0.5)=P(0.2\leq X\leq 0.5)=0.3.
\]
The only difference between the two events is the possible inclusion of the two
endpoints \(0.2\) and \(0.5\), and each endpoint has probability zero.

More generally, for the continuous random variables studied in this chapter,
\[
P(a<X<b)=P(a\leq X\leq b)=P(a<X\leq b)=P(a\leq X<b).
\]
This is convenient, but it should not be used blindly in every probability
model. It relies on the absence of point masses.

\subsection*{What must be learned again}

The discrete chapter introduced random variables, support, probability mass
functions, cumulative distribution functions, expectation, variance, moments, and
classical models. The continuous chapter revisits these ideas, but several
concepts must be reinterpreted:
\begin{itemize}
    \item point probabilities are replaced by interval probabilities;
    \item probability mass functions are replaced by density functions;
    \item sums are replaced by integrals;
    \item probabilities become areas under curves;
    \item expectation becomes a continuous weighted average;
    \item quantiles become especially natural because intervals and cumulative
    probabilities are central.
\end{itemize}

The conceptual difficulty is not the notation itself. The difficulty is that the
same symbols can suggest the wrong discrete intuition. For instance, the symbol
\(f_X(x)\) will describe a density, but \(f_X(x)\) is not usually a probability.
A height of a curve and an area under a curve are different objects.

\begin{theorembox}
The continuous case should not be read as the discrete case with more values.
It is a different way of distributing probability: probability is assigned to
regions by accumulation, not to individual points by mass.
\end{theorembox}
